\begin{abstract}
Modern machine learning (ML) training workloads place substantial demands on both computational and communication resources. Consequently, accurate performance estimation has become increasingly critical for guiding system design decisions, such as the selection of parallelization strategies, cluster configurations, and hardware provisioning. Existing simulation-based performance estimation requires reimplementing the ML framework in a simulator, which demands significant manual effort and is hard to maintain as ML frameworks evolve rapidly.

This paper introduces \sys, a hybrid GPU cluster simulator designed for performance estimation of ML training workloads. \sys executes unmodified ML frameworks as is within a distributed, containerized environment. Each container emulates the behavior of a GPU server in a large-scale cluster, while \sys intercepts and simulates GPU- and communication-related operations to provide high-fidelity performance estimation. We call this approach hybrid simulation of ML systems, in contrast to traditional methods that simulate static workloads. The primary advantage of hybrid simulation is that it allows direct reuse of ML framework source code in simulation, avoiding the need for reimplementation. Our evaluation shows that \sys provides accuracy comparable to static workload simulation while supporting three state-of-the-art LLM training frameworks out-of-the-box. In addition, \sys operates on a single GPU, eliminating the need for the resource-intensive trace collection and workload extraction steps required by traditional trace-based simulators.
\sys is open-sourced at \url{https://github.com/QDelta/Phantora}. 

\end{abstract}